% !TeX root = ../偏微分方程笔记.tex
% 由主文档通过 \input 载入；请编译主文档。
\setcounter{chapter}{7} % 保留原书章号。
\chapter{多维波动方程与能量法}
\label{chap:wave-multi}

高维空间中，波沿球面传播，几何使不同维数的表示公式有所不同。另一方面，把方程与速度相乘可以得到不依赖显式公式的能量关系，这种方法还能推广到变系数和弱解。

\section{球面平均与三维表示公式}
对 $\psi:\mathbb R^3\to\mathbb R$ 定义\emph{球面平均（spherical mean）}
\[
 M_r\psi(x)=\frac1{4\pi}\int_{|\omega|=1}\psi(x+r\omega)\,dS_\omega.
\]
对光滑函数，散度定理给出 $\partial_{rr}M_r\psi+2r^{-1}\partial_rM_r\psi=M_r\Delta\psi$。所以对三维波动解 $u$，函数 $W(r,t)=rM_ru(x,t)$ 满足 $W_{tt}=c^2W_{rr}$，且关于 $r$ 奇对称。对 $W$ 使用一维公式，再由 $u(x,t)=\partial_rW(0,t)$ 恢复解，得到\emph{基尔霍夫公式（Kirchhoff's formula）}
\begin{equation}
 u(x,t)=\partial_t\bigl(tM_{ct}g(x)\bigr)+tM_{ct}h(x),\qquad t>0.
 \label{eq:kirchhoff}
\end{equation}
例如 $g\in C^3$、$h\in C^2$ 是保证这些运算及经典验证成立的一组充分条件。由于 $M_0g=g$ 且球面上一阶方向的平均为零，令 $t\downarrow0$ 可恢复 $u=g$、$u_t=h$。

\section{降维法与传播尾部}
让三维数据不依赖第三个坐标，得到二维解。把单位球面的上下半球投影到单位圆盘，面元之和为 $2(1-|z|^2)^{-1/2}dz$，因此令
\[
 P_t\psi(x)=\frac1{2\pi}\int_{|z|<1}\frac{\psi(x+ctz)}{\sqrt{1-|z|^2}}\,dz,
\]
便有二维公式 $u=\partial_t(tP_tg)+tP_th$。分母虽在边缘奇异，但可积；常数的归一化由 $P_t1=1$ 检查。

三维公式仅使用以 $x$ 为中心、半径 $ct$ 的球面附近的初始数据；二维公式则使用整个圆盘。对于自由空间中的三维常系数波动，紧支集数据产生的信号经过后可以完全消失，体现\emph{惠更斯原理（Huygens' principle）}。二维一般留有传播尾部，不能将球面依赖误推广到所有维数。

\section{能量恒等式与唯一性}
设 $\Omega$ 为有界 $C^1$ 区域，$u$ 足够光滑，满足 $u_{tt}-c^2\Delta u=f$。定义
\begin{equation}
 E(t)=\frac12\int_\Omega\bigl(u_t^2+c^2|\nabla u|^2\bigr)\,dx.
 \label{eq:wave-energy}
\end{equation}
乘方程以 $u_t$ 后分部积分，得到
\begin{equation}
 E'(t)=\int_\Omega f u_t\,dx+c^2\int_{\partial\Omega}u_t\partial_\nu u\,dS.
 \label{eq:wave-energy-balance}
\end{equation}
固定的齐次狄利克雷边界给出 $u_t=0$，齐次诺伊曼边界给出 $\partial_\nu u=0$，两者都消去边界项。无源时 $E(t)=E(0)$。

乘以速度 $u_t$ 的理由是把加速度与速度的乘积写成动能导数：$u_{tt}u_t=\partial_t(u_t^2/2)$；空间分部积分又把 $-\Delta u\,u_t$ 变成 $\nabla u\cdot\nabla u_t=\partial_t(|\nabla u|^2/2)$。因此两种二次量恰好拼成一个全时间导数。对真实弦，乘上线密度后，第一项是动能，第二项与张力产生的弹性势能相应；式\eqref{eq:wave-energy}使用了除以密度后的归一化单位。边界项表示外界通过端点做功，故移动的非零狄利克雷边界一般不能把它删掉。

\begin{example}[能量在两种形式之间交换]
在固定端弦上取 $u(x,t)=A\cos(\omega t)\sin(k\pi x/\ell)$，其中 $\omega=ck\pi/\ell$。利用正弦和余弦平方的积分均为 $\ell/2$，得到
\[
 E_{\mathrm{kin}}(t)=\frac{A^2\omega^2\ell}{4}\sin^2(\omega t),\qquad
 E_{\mathrm{pot}}(t)=\frac{A^2\omega^2\ell}{4}\cos^2(\omega t).
\]
总能量为常数 $A^2\omega^2\ell/4$。最大位移时速度为零、能量全为势能；弦经过平衡位置时势能为零、能量全为动能。能量守恒不表示每一处的位移不变，也不表示动能与势能分别守恒。
\end{example}

\begin{proposition}
对同一组初位移、初速度和上述同类齐次边界，光滑波动解至多有一个。
\end{proposition}
\begin{proof}
两个解之差 $w$ 满足齐次方程、零初值及齐次边界，因而 $E_w(t)=0$，特别是 $w_t=0$。结合 $w(x,0)=0$ 得 $w=0$。这也处理了诺伊曼条件下梯度不能单独控制常数的问题。
\end{proof}
有源时由柯西--施瓦茨不等式得到 $E'\le\sqrt{2E}\|f\|_2$，经对 $\sqrt{E+\varepsilon}$ 积分再令 $\varepsilon\downarrow0$ 得
\[
 \sqrt{E(t)}\le\sqrt{E(0)}+\frac1{\sqrt2}\int_0^t\|f(s)\|_2\,ds.
\]
将此式用于解之差即可得到能量范数中的连续依赖。

\section{局部能量与有限传播}
在无源全空间问题中，固定 $(x_0,t_0)$，对缩小的球 $B_{c(t_0-t)}(x_0)$ 定义局部能量。求导时球面向内移动，除式\eqref{eq:wave-energy-balance}的通量外，还产生负的运动边界项，因此
\[
 E_{\mathrm{loc}}'(t)=\int_{\partial B_{c(t_0-t)}}
 \left[c^2u_t\partial_\nu u-\frac c2(u_t^2+c^2|\nabla u|^2)\right]dS\le0.
\]
最后一步来自 $2c|u_t\partial_\nu u|\le u_t^2+c^2|\nabla u|^2$。若两组数据在初始球内相同，差解的局部能量为零，结合初位移便得到过去光锥中的差解为零。这给出任意空间维数的有限传播证明。

阻尼方程 $u_{tt}+\gamma u_t-c^2\Delta u=0$ 在零通量边界下满足 $E'=-\gamma\int u_t^2\le0$。该式直接给出单调性；指数衰减还需把速度项与总能量联系起来，不能仅由 $E'\le0$ 推出。

\section{例题与习题}
\begin{example}
三维径向解 $u(x,t)=U(r,t)$，$r=|x|$，如何化为一维问题？
\end{example}
\begin{solve}
对 $r>0$，$\Delta u=U_{rr}+2U_r/r$。令 $v=rU$，则 $v_{tt}=c^2v_{rr}$。若 $U$ 在原点光滑，则 $v(0,t)=0$，应把 $v$ 奇延拓后使用一维公式。求出 $v$ 再除以 $r$；原点值用极限 $U(0,t)=v_r(0,t)$ 定义，不能直接代入含 $1/r$ 的表达式。
\end{solve}
\begin{enumerate}
 \item 验证 $M_r1=1$，用式\eqref{eq:kirchhoff}处理常数初位移与常数初速度。
 \item 对齐次狄利克雷问题证明能量估计中 $\|\nabla u\|_2$ 能控制 $\|u\|_2$。
 \item 写出含阻尼和源项时的完整能量收支，解释各项符号。
\end{enumerate}

